% --- Insert Image ------------------------------------------------------------
% Usage: \insertimage{filename}{caption}{label}{width}
% If the image file does not exist yet, a shaded placeholder is rendered
% so the document builds cleanly while diagrams are pending.
\newcommand{\insertimage}[4]{%
  \begin{figure}[H]
    \centering
    \IfFileExists{#1}{%
      \includegraphics[width=#4]{#1}%
    }{%
      \fcolorbox{gray}{gray!15}{%
        \parbox[c][6cm]{#4}{%
          \centering\color{gray}%
          \textit{[Pending: \texttt{#1}]}%
        }%
      }%
    }%
    \caption{#2}
    \label{#3}
  \end{figure}
}

% --- Insert Two Images Side by Side -----------------------------------------
% Usage: \inserttwoimages{file1}{cap1}{label1}{file2}{cap2}{label2}{overall_caption}{overall_label}
\newcommand{\inserttwoimages}[8]{%
  \begin{figure}[H]
    \centering
    \begin{subfigure}[b]{0.48\textwidth}
      \centering
      \includegraphics[width=\textwidth]{#1}
      \caption{#2}
      \label{#3}
    \end{subfigure}
    \hfill
    \begin{subfigure}[b]{0.48\textwidth}
      \centering
      \includegraphics[width=\textwidth]{#4}
      \caption{#5}
      \label{#6}
    \end{subfigure}
    \caption{#7}
    \label{#8}
  \end{figure}
}

% --- Insert Two Mobile Screenshots Side by Side -----------------------------
% Same interface as \inserttwoimages but with narrower subfigures (0.38)
% for mobile application screens that are taller than wide.
\newcommand{\inserttwomobileimages}[8]{%
  \begin{figure}[H]
    \centering
    \begin{subfigure}[b]{0.35\textwidth}
      \centering
      \includegraphics[width=\textwidth]{#1}
      \caption{#2}
      \label{#3}
    \end{subfigure}
    \hfill
    \begin{subfigure}[b]{0.35\textwidth}
      \centering
      \includegraphics[width=\textwidth]{#4}
      \caption{#5}
      \label{#6}
    \end{subfigure}
    \caption{#7}
    \label{#8}
  \end{figure}
}

% --- Insert Table ------------------------------------------------------------
% Usage: \inserttable{caption}{label}{column_spec}{rows}
%   column_spec : any tabularx column spec, e.g. {l X r} or {lcc}
%   rows        : table body including \toprule … \bottomrule
% The macro uses tabularx at \textwidth and sets arraystretch to 1.25.
\newcommand{\inserttable}[4]{%
  \begin{table}[H]
    \centering
    \renewcommand{\arraystretch}{1.25}%
    \begin{tabularx}{\textwidth}{#3}
      #4
    \end{tabularx}
    \caption{#1}
    \label{#2}
  \end{table}
}

% --- Insert Code from External File -----------------------------------------
% Usage: \insertcode{filepath}{language}{caption}{label}
%   e.g. \insertcode{code/sample.py}{python}{Sample script}{lst:sample}
\newcommand{\insertcode}[4]{%
  \begin{listing}[H]
    \inputminted[linenos, frame=single, breaklines, label=#1]{#2}{#1}
    \caption{#3}
    \label{#4}
  \end{listing}
}

% --- Insert Specific Lines from External File -------------------------------
% Usage: \insertcodelines{filepath}{language}{firstline}{lastline}{caption}{label}
%   e.g. \insertcodelines{code/sample.py}{python}{5}{15}{Main function}{lst:main}
\newcommand{\insertcodelines}[6]{%
  \begin{listing}[H]
    \inputminted[linenos, frame=single, breaklines, firstline=#3, lastline=#4, label=#1]{#2}{#1}
    \caption{#5}
    \label{#6}
  \end{listing}
}

% --- Insert Code from File (NO filename shown) ------------------------------
% Like \insertcode but the filename is NOT displayed on the code block.
% Usage: \insertcodeclean{filepath}{language}{caption}{label}
%   e.g. \insertcodeclean{code/sample.py}{python}{API handler}{lst:handler}
\newcommand{\insertcodeclean}[4]{%
  \begin{listing}[H]
    \inputminted[linenos, frame=single, breaklines, fontsize=\small]{#2}{#1}
    \caption{#3}
    \label{#4}
  \end{listing}
}

% --- Insert Code lines from File (NO filename shown) ------------------------------
% Like \insertcode but the filename is NOT displayed on the code block.
% Usage: \insertlinescodeclean{filepath}{language}{first}{last}{caption}{label}
%   e.g. \insertlinescodeclean{code/sample.py}{python}{4}{12}{API handler}{lst:handler}
\newcommand{\insertlinescodeclean}[6]{%
  \begin{listing}[H]
    \inputminted[linenos, frame=single, breaklines, firstline=#3, lastline=#4, fontsize=\small]{#2}{#1}
    \caption{#5}
    \label{#6}
  \end{listing}
}

% --- Insert Code Block from File (NO filename, NO caption) ------------------
% Like \codeblock but reads from a file instead of inline.
% Usage: \insertcodeblock{filepath}{language}
%   e.g. \insertcodeblock{code/sample.py}{python}
\newcommand{\insertcodeblock}[2]{%
  \inputminted[linenos, frame=single, breaklines, fontsize=\small]{#2}{#1}
}

% --- Insert Reference (shorthand) -------------------------------------------
% Usage: \insertref{key}
\newcommand{\insertref}[1]{\cite{#1}}

% --- Note / Callout Box -----------------------------------------------------
% Usage: \note{This is an important note.}
\newtcolorbox{notebox}{
  colback=blue!5!white,
  colframe=NavyBlue!70!black,
  fonttitle=\bfseries,
  title=Note,
  arc=3pt,
  boxrule=0.8pt,
  left=8pt, right=8pt, top=6pt, bottom=6pt
}
\newcommand{\note}[1]{\begin{notebox}#1\end{notebox}}

% --- Warning Box ------------------------------------------------------------
% Usage: \warningbox{Be careful with this step.}
\newtcolorbox{warnbox}{
  colback=red!5!white,
  colframe=BrickRed!70!black,
  fonttitle=\bfseries,
  title=Warning,
  arc=3pt,
  boxrule=0.8pt,
  left=8pt, right=8pt, top=6pt, bottom=6pt
}
\newcommand{\warningbox}[1]{\begin{warnbox}#1\end{warnbox}}

% --- Tip Box ----------------------------------------------------------------
% Usage: \tipbox{A useful tip for the reader.}
\newtcolorbox{tipboxenv}{
  colback=green!5!white,
  colframe=ForestGreen!70!black,
  fonttitle=\bfseries,
  title=Tip,
  arc=3pt,
  boxrule=0.8pt,
  left=8pt, right=8pt, top=6pt, bottom=6pt
}
\newcommand{\tipbox}[1]{\begin{tipboxenv}#1\end{tipboxenv}}

% --- Inline Code Block (no external file) -----------------------------------
% Usage: \begin{inlinecode}{language}{caption}{label}
%          your code here...
%        \end{inlinecode}
\makeatletter
\newcommand{\@inlinecode@caption}{}
\newcommand{\@inlinecode@label}{}
\newenvironment{inlinecode}[3]{%
  \renewcommand{\@inlinecode@caption}{#2}%
  \renewcommand{\@inlinecode@label}{#3}%
  \VerbatimEnvironment
  \begin{listing}[H]%
  \begin{minted}[linenos, frame=single, breaklines, fontsize=\small]{#1}%
}{%
  \end{minted}%
  \caption{\@inlinecode@caption}%
  \label{\@inlinecode@label}%
  \end{listing}%
}
\makeatother

% Simple inline code (no caption/label, just the block)
% Usage: \begin{codeblock}{language}
%          code here...
%        \end{codeblock}
%        \end{codeblock}
\newenvironment{codeblock}[1]{%
  \VerbatimEnvironment
  \begin{minted}[linenos, frame=single, breaklines, fontsize=\small]{#1}%
}{%
  \end{minted}%
}

% --- Code Diff (uses minted with Pygments diff lexer) -----------------------
% Auto-generate .diff files with: python3 scripts/codediff.py <old> <new> -o <out.diff>

% Insert a diff from an external .diff file (with caption & label)
% Usage: \insertdiff{filepath}{caption}{label}
%   e.g. \insertdiff{code/refactor.diff}{Refactoring to async}{lst:diff}
\newcommand{\insertdiff}[3]{%
  \begin{listing}[H]
    \inputminted[frame=single, breaklines, fontsize=\small]{diff}{#1}
    \caption{#2}
    \label{#3}
  \end{listing}
}

% Write diff inline (no external file, with caption & label)
% Usage: \begin{codediff}{caption}{label}
%          --- old.py
%          +++ new.py
%          -removed line
%          +added line
%        \end{codediff}
\makeatletter
\newcommand{\@codediff@caption}{}
\newcommand{\@codediff@label}{}
\newenvironment{codediff}[2]{%
  \renewcommand{\@codediff@caption}{#1}%
  \renewcommand{\@codediff@label}{#2}%
  \VerbatimEnvironment
  \begin{listing}[H]%
  \begin{minted}[frame=single, breaklines, fontsize=\small]{diff}%
}{%
  \end{minted}%
  \caption{\@codediff@caption}%
  \label{\@codediff@label}%
  \end{listing}%
}
\makeatother

% --- Compare Two Implementations Side by Side --------------------------------
% Usage: \comparecode{title1}{lang1}{code1}{title2}{lang2}{code2}{caption}{label}
\newcommand{\comparecode}[8]{%
  \begin{figure}[H]
    \centering
    \begin{minipage}[t]{0.48\textwidth}
      \centering
      \textbf{#1}
      \inputminted[linenos, frame=single, breaklines, fontsize=\scriptsize]{#2}{#3}
    \end{minipage}
    \hfill
    \begin{minipage}[t]{0.48\textwidth}
      \centering
      \textbf{#4}
      \inputminted[linenos, frame=single, breaklines, fontsize=\scriptsize]{#5}{#6}
    \end{minipage}
    \caption{#7}
    \label{#8}
  \end{figure}
}
